\documentclass{article}

\usepackage[utf8]{inputenc}
\usepackage[spanish]{babel}
\usepackage{tikz}

% Cargamos las bibliotecas necesarias
\usetikzlibrary{shapes,arrows,positioning,calc}

% --- Estilos locales para el diagrama TikZ ---
\tikzset{
	block/.style = {
		draw, 
		fill=blue!10, 
		rectangle, 
		minimum height=3em, 
		minimum width=5em, 
		line width=1pt, 
		align=center, 
		rounded corners=2pt 
	},
	sum/.style = {
		draw, 
		fill=blue!10, 
		circle, 
		minimum size=1.5em, 
		line width=1pt,
		inner sep=0pt
	},
	signal/.style = {
		->, 
		>=stealth', 
		thick, 
		shorten >=1pt, 
		shorten <=1pt
	},
	label/.style = {
		font=\small, 
		midway, 
		above,
		align=center
	},
	branch/.style = {
		fill, 
		shape=circle, 
		minimum size=4pt, 
		inner sep=0pt
	}
}
% --------------------------------------------

\begin{document}
	
	\begin{figure}[htbp]
		\centering
		\begin{tikzpicture}[auto]
			
			% --- Definición de NODOS ---
			% Nodo de entrada invisible
			\node (input) at (0,0) {};
			
			% Sumador (donde se mezclan la entrada y la realimentación)
			\node [sum] (sum) at (2.5,0) {$+$};
			
			% Sistema 1 (Ruta de avance / Forward path)
			\node [block] (sys1) at (5,0) {Sistema 1 \\ $S_1$};
			
			% Punto de bifurcación para tomar la muestra de salida
			\node [branch] (branch) at (7.5,0) {};
			
			% Nodo de salida invisible
			\node (output) at (9.5,0) {};
			
			% Sistema 2 (Ruta de realimentación / Feedback path)
			% Lo colocamos debajo del Sistema 1
			\node [block] (sys2) at (5,-2) {Sistema 2 \\ $S_2$};
			
			% --- CONEXIONES ---
			% Entrada -> Sumador (Añadimos un pequeño "+" para indicar suma positiva)
			\draw [signal] (input) -- node[label] {Entrada, $x(t)$} 
			node[pos=0.85, above left, font=\footnotesize] {$+$} (sum);
			
			% Sumador -> Sistema 1 (Señal de error)
			\draw [signal] (sum) -- node[label] {$e(t)$} (sys1);
			
			% Sistema 1 -> Bifurcación -> Salida
			\draw [signal] (sys1) -- (branch) -- node[label] {Salida, $y(t)$} (output);
			
			% Bifurcación -> Sistema 2 (Baja y dobla a la izquierda)
			\draw [signal] (branch) |- (sys2);
			
			% Sistema 2 -> Sumador (Sigue a la izquierda y sube al sumador)
			% Añadimos un pequeño "-" para indicar realimentación negativa (típico en control)
			\draw [signal] (sys2) -| node[pos=0.9, right, font=\footnotesize] {$-$} (sum.south);
			
			% --- CAJA PUNTEADA ---
			% Coordenadas absolutas que envuelven todo el lazo (desde antes del sumador hasta la bifurcación)
			\draw[dashed, gray, thick, rounded corners=5pt]
			(1.5, 1.2) rectangle (8.2, -3.2);
			
			% Etiqueta inferior de la caja
			\node[font=\small] at (4.85, -2.9) {Sistema Total Comprendido};
			
		\end{tikzpicture}
		\caption{Ejemplo de interconexión de sistemas con realimentación (feedback). La salida del Sistema 1 se mide y se hace pasar por el Sistema 2, para luego restarse a la señal de entrada original.}
		\label{fig:interconexion_realimentacion}
	\end{figure}
	
\end{document}